% GENERATED FILE. Edit canonical lessons, then run npm run generate:books.
% canonical-source-hash: fnv1a64:bea54141df36961d
% canonical-lessons: KA-C520-todagu, KA-C520-vahisu, KA-C520-palagu, KA-C520-melvicarane, KA-C520-samayapalane

\chapter{Teamwork}
\label{ch:ka-teamwork}
\hlchaptermodality{520}

\begin{chapteropening}
\textbf{By the end of this chapter:} \emph{I can say to get involved, to take charge of, to get practised, supervision and punctuality.}

Complete the last lesson of chapter 520: I can say to get involved, to take charge of, to get practised, supervision and punctuality.
\end{chapteropening}

\section[\texorpdfstring{\kn{ತೊಡಗು} (toḍagu)}{ತೊಡಗು (toḍagu)}]{\kn{ತೊಡಗು} (toḍagu) — to get involved; to set to work}

\label{lesson:KA-C520-todagu}



\begin{quote}
Before the new one: say the Kannada for to assign, then the Kannada for to save up.
\end{quote}

\subsection*{You'll want to know: \kn{ತೊಡಗು}}
\textbf{\kn{ತೊಡಗು}} — \emph{toḍagu} — \textquotedblleft{}to get involved; to set to work\textquotedblright{}.

\begin{quote}
\kn{ಊರಿನ} \kn{ಯುವಕರು} \kn{ಕೆರೆಯ} \kn{ಕೆಲಸದಲ್ಲಿ} \kn{ತೊಡಗಿದರು}.
\end{quote}

\begin{quote}
\emph{ūrina yuvakaru kereya kelasadalli toḍagidaru.}
\end{quote}

\begin{quote}
The young men of the town set to work on the tank.
\end{quote}

Past (I …): \kn{ತೊಡಗಿದೆ}, \emph{toḍagide}.

In what: -\kn{ಅಲ್ಲಿ}. \kn{ತೊಡಗಿಸಿಕೊ}, to get yourself involved.

\begin{cousinweb}
\textbf{\kn{ತೊಡಗು}} (\emph{toḍagu}) is native Dravidian: DEDR 3481 lists Kannada \emph{toḍagu}, \textquotedblleft{}to engage (as in work), begin, undertake\textquotedblright{}.
\end{cousinweb}

\begin{grammarlens}[title={What you've built}]
One more word for describing what you lived and felt.
\end{grammarlens}

\subsection*{Your turn}
\begin{itemize}
\raggedright
  \item \emph{Say it:} \emph{toḍagu}
  \item \emph{Say it:} \emph{toḍagu}, once more
  \item \emph{Say it:} \emph{ūrina yuvakaru kereya kelasadalli toḍagidaru.} — The young men of the town set to work on the tank.
  \item \emph{Recall:} say the Kannada for to assign, then the Kannada for to save up, then say \emph{toḍagu} again
\end{itemize}

\begin{tcolorbox}[breakable,colback=teal!4,colframe=teal!35!black,title={Before you move on}]
What does \kn{ತೊಡಗು} mean? (\textquotedblleft{}To get involved; to set to work\textquotedblright{}.) Say it once more.
\end{tcolorbox}
\section[\texorpdfstring{\kn{ವಹಿಸು} (vahisu)}{ವಹಿಸು (vahisu)}]{\kn{ವಹಿಸು} (vahisu) — to take charge of; to take on}

\label{lesson:KA-C520-vahisu}



\begin{quote}
Before the new one: say the Kannada for to save up, then the Kannada for to get involved.
\end{quote}

\subsection*{You'll want to know: \kn{ವಹಿಸು}}
\textbf{\kn{ವಹಿಸು}} — \emph{vahisu} — \textquotedblleft{}to take charge of; to take on\textquotedblright{}.

\begin{quote}
\kn{ಹಬ್ಬದ} \kn{ಊಟದ} \kn{ಜವಾಬ್ದಾರಿಯನ್ನು} \kn{ಮಾವ} \kn{ವಹಿಸಿದರು}.
\end{quote}

\begin{quote}
\emph{habbada ūṭada javābdāriyannu māva vahisidaru.}
\end{quote}

\begin{quote}
Uncle took charge of the festival meal.
\end{quote}

Past (I …): \kn{ವಹಿಸಿದೆ}, \emph{vahiside}.

Often with \kn{ಜವಾಬ್ದಾರಿ} or \kn{ಕಾಳಜಿ}: \kn{ಕಾಳಜಿ} \kn{ವಹಿಸು}, to take care.

\begin{cousinweb}
\textbf{\kn{ವಹಿಸು}} (\emph{vahisu}) is built on the Sanskrit root \emph{vah}, \textquotedblleft{}to carry, to bear\textquotedblright{} (Monier-Williams), with the verb ending -\emph{isu}: to shoulder a task.
\end{cousinweb}

\begin{grammarlens}[title={What you've built}]
One more word for describing what you lived and felt.
\end{grammarlens}

\subsection*{Your turn}
\begin{itemize}
\raggedright
  \item \emph{Say it:} \emph{vahisu}
  \item \emph{Say it:} \emph{vahisu}, once more
  \item \emph{Say it:} \emph{habbada ūṭada javābdāriyannu māva vahisidaru.} — Uncle took charge of the festival meal.
  \item \emph{Recall:} say the Kannada for to save up, then the Kannada for to get involved, then say \emph{vahisu} again
\end{itemize}

\begin{tcolorbox}[breakable,colback=teal!4,colframe=teal!35!black,title={Before you move on}]
What does \kn{ವಹಿಸು} mean? (\textquotedblleft{}To take charge of; to take on\textquotedblright{}.) Say it once more.
\end{tcolorbox}
\section[\texorpdfstring{\kn{ಪಳಗು} (paḷagu)}{ಪಳಗು (paḷagu)}]{\kn{ಪಳಗು} (paḷagu) — to get practised; to become skilled}

\label{lesson:KA-C520-palagu}



\begin{quote}
Before the new one: say the Kannada for to get involved, then the Kannada for to take charge of.
\end{quote}

\subsection*{You'll want to know: \kn{ಪಳಗು}}
\textbf{\kn{ಪಳಗು}} — \emph{paḷagu} — \textquotedblleft{}to get practised; to become skilled\textquotedblright{}.

\begin{quote}
\kn{ಎರಡು} \kn{ವರ್ಷದಲ್ಲಿ} \kn{ಅವನು} \kn{ಈ} \kn{ಕೆಲಸದಲ್ಲಿ} \kn{ಚೆನ್ನಾಗಿ} \kn{ಪಳಗಿದನು}.
\end{quote}

\begin{quote}
\emph{eraḍu varṣadalli avanu ī kelasadalli cennāgi paḷagidanu.}
\end{quote}

\begin{quote}
In two years he became well practised in this work.
\end{quote}

Past (I …): \kn{ಪಳಗಿದೆ}, \emph{paḷagide}.

In what: -\kn{ಅಲ್ಲಿ}. \kn{ಪಳಗಿಸು}, to tame or train an animal.

\begin{cousinweb}
\textbf{\kn{ಪಳಗು}} (\emph{paḷagu}) is native Dravidian: DEDR 4000 lists older Kannada \emph{paḻagu}, \textquotedblleft{}to become used to, to be trained\textquotedblright{}, and \emph{paḻaka}, \textquotedblleft{}use, practice, habit\textquotedblright{}, beside Tamil \emph{paḻaku}.
\end{cousinweb}

\begin{grammarlens}[title={What you've built}]
One more word for describing what you lived and felt.
\end{grammarlens}

\subsection*{Your turn}
\begin{itemize}
\raggedright
  \item \emph{Say it:} \emph{paḷagu}
  \item \emph{Say it:} \emph{paḷagu}, once more
  \item \emph{Say it:} \emph{eraḍu varṣadalli avanu ī kelasadalli cennāgi paḷagidanu.} — In two years he became well practised in this work.
  \item \emph{Recall:} say the Kannada for to get involved, then the Kannada for to take charge of, then say \emph{paḷagu} again
\end{itemize}

\begin{tcolorbox}[breakable,colback=teal!4,colframe=teal!35!black,title={Before you move on}]
What does \kn{ಪಳಗು} mean? (\textquotedblleft{}To get practised; to become skilled\textquotedblright{}.) Say it once more.
\end{tcolorbox}
\section[\texorpdfstring{\kn{ಮೇಲ್ವಿಚಾರಣೆ} (mēlvicāraṇe)}{ಮೇಲ್ವಿಚಾರಣೆ (mēlvicāraṇe)}]{\kn{ಮೇಲ್ವಿಚಾರಣೆ} (mēlvicāraṇe) — supervision}

\label{lesson:KA-C520-melvicarane}



\begin{quote}
Before the new one: say the Kannada for to take charge of, then the Kannada for to get practised.
\end{quote}

\subsection*{You'll want to know: \kn{ಮೇಲ್ವಿಚಾರಣೆ}}
\textbf{\kn{ಮೇಲ್ವಿಚಾರಣೆ}} — \emph{mēlvicāraṇe} — \textquotedblleft{}supervision\textquotedblright{}.

\begin{quote}
\kn{ಎಂಜಿನಿಯರ್} \kn{ಮೇಲ್ವಿಚಾರಣೆಯಲ್ಲಿ} \kn{ಕೆಲಸ} \kn{ನಡೆಯಿತು}.
\end{quote}

\begin{quote}
\emph{eñjiniyar mēlvicāraṇeyalli kelasa naḍeyitu.}
\end{quote}

\begin{quote}
The work went on under the engineer's supervision.
\end{quote}

\kn{ಮೇಲ್}, over + \kn{ವಿಚಾರಣೆ}, looking into: overseeing.

\begin{cousinweb}
\textbf{\kn{ಮೇಲ್ವಿಚಾರಣೆ}} (\emph{mēlvicāraṇe}) joins native \emph{mēl}, \textquotedblleft{}above, over\textquotedblright{} (DEDR 5086), and Sanskrit \emph{vicāraṇa}, \textquotedblleft{}consideration, examination\textquotedblright{} (Monier-Williams).
\end{cousinweb}

\begin{grammarlens}[title={What you've built}]
One more word for describing what you lived and felt.
\end{grammarlens}

\subsection*{Your turn}
\begin{itemize}
\raggedright
  \item \emph{Say it:} \emph{mēlvicāraṇe}
  \item \emph{Say it:} \emph{mēlvicāraṇe}, once more
  \item \emph{Say it:} \emph{eñjiniyar mēlvicāraṇeyalli kelasa naḍeyitu.} — The work went on under the engineer's supervision.
  \item \emph{Recall:} say the Kannada for to take charge of, then the Kannada for to get practised, then say \emph{mēlvicāraṇe} again
\end{itemize}

\begin{tcolorbox}[breakable,colback=teal!4,colframe=teal!35!black,title={Before you move on}]
What does \kn{ಮೇಲ್ವಿಚಾರಣೆ} mean? (\textquotedblleft{}Supervision\textquotedblright{}.) Say it once more.
\end{tcolorbox}
\section[\texorpdfstring{\kn{ಸಮಯಪಾಲನೆ} (samayapālane)}{ಸಮಯಪಾಲನೆ (samayapālane)}]{\kn{ಸಮಯಪಾಲನೆ} (samayapālane) — punctuality}

\label{lesson:KA-C520-samayapalane}



\begin{quote}
Before the new one: say the Kannada for to get practised, then the Kannada for supervision.
\end{quote}

\subsection*{You'll want to know: \kn{ಸಮಯಪಾಲನೆ}}
\textbf{\kn{ಸಮಯಪಾಲನೆ}} — \emph{samayapālane} — \textquotedblleft{}punctuality\textquotedblright{}.

\begin{quote}
\kn{ಸಮಯಪಾಲನೆ} \kn{ಒಳ್ಳೆಯ} \kn{ನೌಕರನ} \kn{ಗುಣ}.
\end{quote}

\begin{quote}
\emph{samayapālane oḷḷeya naukarana guṇa.}
\end{quote}

\begin{quote}
Punctuality is the mark of a good employee.
\end{quote}

\kn{ಸಮಯ}, time + \kn{ಪಾಲನೆ}, keeping.

\begin{cousinweb}
\textbf{\kn{ಸಮಯಪಾಲನೆ}} (\emph{samayapālane}) is a learned compound built from Sanskrit \emph{samaya}, \textquotedblleft{}a meeting, an agreed time\textquotedblright{} (Monier-Williams), and \emph{pālana}, \textquotedblleft{}protecting, keeping\textquotedblright{} (CDIAL 8128).
\end{cousinweb}

\begin{grammarlens}[title={What you've built}]
One more word for describing what you lived and felt.
\end{grammarlens}

\subsection*{Your turn}
\begin{itemize}
\raggedright
  \item \emph{Say it:} \emph{samayapālane}
  \item \emph{Say it:} \emph{samayapālane}, once more
  \item \emph{Say it:} \emph{samayapālane oḷḷeya naukarana guṇa.} — Punctuality is the mark of a good employee.
  \item \emph{Recall:} say the Kannada for to get practised, then the Kannada for supervision, then say \emph{samayapālane} again
\end{itemize}

\begin{tcolorbox}[breakable,colback=teal!4,colframe=teal!35!black,title={Before you move on}]
What does \kn{ಸಮಯಪಾಲನೆ} mean? (\textquotedblleft{}Punctuality\textquotedblright{}.) Say it once more.
\end{tcolorbox}
